\section*{Chapter \ref{ch:g8:plastics} --- Plastics and Synthetic Materials}

\begin{solution}{exo:g8:plastics:1}
Natural: cotton, wool. Artificial: paper, viscose. Synthetic: nylon,
polystyrene.
\end{solution}

\begin{solution}{exo:g8:plastics:2}
From crude oil (or natural gas).
\end{solution}

\begin{solution}{exo:g8:plastics:3}
PET (polyethylene terephthalate): water and soda bottles.
\end{solution}

\begin{solution}{exo:g8:plastics:4}
A \omterm{def:g8:plastics:thermoplastic}{thermoplastic} softens each time it is heated and can be remoulded; a
\omterm{def:g8:plastics:thermoplastic}{thermoset} sets for good the first time and chars instead of softening.
\end{solution}

\begin{solution}{exo:g8:plastics:5}
Polypropylene (PP), a \omterm{def:g8:plastics:thermoplastic}{thermoplastic}: it can be melted and remoulded.
\end{solution}

\begin{solution}{exo:g8:plastics:6}
The triangle only names the \omterm{def:g8:plastics:plastic}{plastic}, to help sorting. Whether the object
is recycled depends on its being collected and on a plant that recycles
that \omterm{def:g8:plastics:plastic}{plastic}.
\end{solution}

\begin{solution}{exo:g8:plastics:7}
A pan gets hot: a \omterm{def:g8:plastics:thermoplastic}{thermoplastic} would soften, a \omterm{def:g8:plastics:thermoplastic}{thermoset} does not.
\end{solution}

\begin{solution}{exo:g8:plastics:8}
$\frac{353}{460} \approx 0.77$: about \qty{77}{\percent}.
\end{solution}

\begin{solution}{exo:g8:plastics:9}
PVC contains chlorine: \omterm{def:g4:what-a-fire-needs:burning}{burning} it gives off hydrogen chloride, a
corrosive gas, besides soot and carbon monoxide from the \omterm{def:g8:combustion-and-fuels:complete}{incomplete
combustion} of a garden fire.
\end{solution}

\begin{solution}{exo:g8:plastics:10}
$300 \times 36 = 10\,800$ bottles; $10\,800 \times 25 = 270\,000$\,g, that
is \qty{270}{kg}.
\end{solution}

\begin{solution}{exo:g8:plastics:11}
Carbon dioxide and water. Its carbon comes from crude oil, stored
underground for millions of years: \omterm{def:g4:what-a-fire-needs:burning}{burning} it adds carbon dioxide to the
\omterm{def:g4:air-a-mixture-of-gases:air}{air}, as \omterm{def:g4:what-a-fire-needs:burning}{burning} a \omterm{def:g8:combustion-and-fuels:fossil}{fossil fuel} does.
\end{solution}

\begin{solution}{exo:g8:plastics:12}
The bottle does not rot: sun and waves break it into smaller and smaller
pieces. The tiny pieces float or sink in the water, and fish swallow
them with their food.
\end{solution}

\begin{solution}{pb:g8:plastics:1}
\textbf{1.} Bodies: 1 (PET). Caps: 2 (HDPE) or 5 (PP).

\textbf{2.} Synthetic: they are made from substances built by chemists,
mostly from oil.

\textbf{3.} Yes: PET is a \omterm{def:g8:plastics:thermoplastic}{thermoplastic}.

\textbf{4.} Each \omterm{def:g8:plastics:plastic}{plastic} is recycled on its own: melted together, a
\omterm{def:g6:pure-substances-and-mixtures:mixture}{mixture} of \omterm{def:g8:plastics:plastic}{plastics} gives a poor material.

\textbf{5.} $0.80 \times 1000 = \qty{800}{kg}$.

\textbf{6.} Caps: $0.12 \times 1000 = \qty{120}{kg}$. Labels: $0.05
\times 1000 = \qty{50}{kg}$.

\textbf{7.} $0.03 \times 1000 = \qty{30}{kg}$ of dirt and liquid.

\textbf{8.} $80 + 12 + 5 = \qty{97}{\percent}$.

\textbf{9.} $0.025 \times 800 = \qty{20}{kg}$.

\textbf{10.} The flakes save crude oil and the energy of making new PET,
and they turn waste into a useful object.

\textbf{11.} Carbon dioxide (and, if the \omterm{def:g4:what-a-fire-needs:burning}{burning} is poor, carbon
monoxide and soot).

\textbf{12.} $800 - 20 = \qty{780}{kg}$ of clean PET flakes per bale.
\end{solution}
